\documentclass[11pt,a4paper]{article}
\usepackage[T1]{fontenc}
\usepackage{amsmath,amssymb,amsthm,mathtools}
\usepackage{geometry}
\usepackage[hidelinks]{hyperref}
\geometry{margin=1in}

\newtheorem{theorem}{Theorem}[section]
\newtheorem{lemma}[theorem]{Lemma}
\newtheorem{proposition}[theorem]{Proposition}
\newtheorem{corollary}[theorem]{Corollary}
\theoremstyle{definition}
\newtheorem{definition}[theorem]{Definition}
\newtheorem{example}[theorem]{Example}

\author{Part II}
\date{Recorded from the lecturer in 2026-2027.}

\title{Stochastic Financial Models}
\begin{document}
\maketitle
\thispagestyle{empty}
\section{Martingales and Conditional Expectation}

Throughout, we work on a probability space
$(\Omega,\mathcal{F},\mathbb{P})$.
All random variables are real-valued unless otherwise stated.

\subsection{Probability spaces and sigma-algebras}

\paragraph{Definition: sigma-algebra.}
A collection $\mathcal{F}\subseteq\mathcal{P}(\Omega)$ is a
\emph{$\sigma$-algebra} on $\Omega$ if:
\begin{enumerate}
    \item $\Omega\in\mathcal{F}$;
    \item if $A\in\mathcal{F}$, then $A^c\in\mathcal{F}$;
    \item if $A_1,A_2,\ldots\in\mathcal{F}$, then
    \[
    \bigcup_{k=1}^{\infty}A_k\in\mathcal{F}.
    \]
\end{enumerate}
These properties also imply that
\[
\varnothing\in\mathcal{F},
\qquad
\bigcap_{k=1}^{\infty}A_k\in\mathcal{F}
\quad\text{whenever }A_k\in\mathcal{F}.
\]

Here $\mathcal{P}(\Omega)$ denotes the power set of $\Omega$.
Thus, $\mathcal{F}$ is a collection of subsets of $\Omega$,
and each $A\in\mathcal{F}$ is called an \emph{event}.

A probability space consists of:
\begin{itemize}
    \item a sample space $\Omega$;
    \item a $\sigma$-algebra $\mathcal{F}$ of events;
    \item a countably additive probability measure
    $\mathbb{P}:\mathcal{F}\to[0,1]$ with
    $\mathbb{P}(\Omega)=1$.
\end{itemize}

\paragraph{Generated sigma-algebra.}
For a collection $\mathcal{C}\subseteq\mathcal{P}(\Omega)$,
the notation $\sigma(\mathcal{C})$ denotes the smallest
$\sigma$-algebra containing $\mathcal{C}$.
It is the intersection of all $\sigma$-algebras on $\Omega$
that contain $\mathcal{C}$.

\paragraph{Almost sure equality.}
Two random variables $U$ and $V$ are equal
\emph{almost surely}, written $U=V$ a.s., if
\[
\mathbb{P}\bigl(\{\omega\in\Omega:U(\omega)=V(\omega)\}\bigr)=1.
\]
Thus, they may differ on a set of probability zero.

\subsection{Borel sets and measurability}

\paragraph{Definition: Borel sigma-algebra.}
The \emph{Borel $\sigma$-algebra} on $\mathbb{R}^d$ is
\[
\mathcal{B}(\mathbb{R}^d)
=
\sigma\bigl(
\{U\subseteq\mathbb{R}^d:U\text{ is open}\}
\bigr).
\]
It is the smallest $\sigma$-algebra containing every open
subset of $\mathbb{R}^d$.

An element $B\in\mathcal{B}(\mathbb{R}^d)$ is called a
\emph{Borel set}.
Borel sets include open sets, closed sets, singletons,
countable sets, and countable unions or intersections of
Borel sets.

For example,
\[
(a,b),\qquad
(-\infty,a],\qquad
\{a\},\qquad
\mathbb{Q}
\]
are Borel subsets of $\mathbb{R}$.
Moreover,
\[
\mathcal{B}(\mathbb{R})
=
\sigma\bigl(
\{(-\infty,a]:a\in\mathbb{R}\}
\bigr).
\]

The notation distinguishes a single set from a collection:
\[
B\subseteq\mathbb{R}^d,
\qquad
B\in\mathcal{B}(\mathbb{R}^d),
\qquad
\mathcal{B}(\mathbb{R}^d)
\subseteq\mathcal{P}(\mathbb{R}^d).
\]

\paragraph{Definition: measurable random variable.}
Let $\mathcal{G}\subseteq\mathcal{F}$ be a $\sigma$-algebra.
A function $X:\Omega\to\mathbb{R}$ is
\emph{$\mathcal{G}$-measurable} if
\[
X^{-1}(B)
=
\{\omega\in\Omega:X(\omega)\in B\}
\in\mathcal{G}
\qquad
\text{for every }B\in\mathcal{B}(\mathbb{R}).
\]
The expression $X^{-1}(B)$ denotes a preimage; it does not
require $X$ to have an inverse function.

Equivalently, it is enough to check that
\[
\{X\leq a\}\in\mathcal{G}
\qquad
\text{for every }a\in\mathbb{R}.
\]
A real-valued \emph{random variable} on
$(\Omega,\mathcal{F},\mathbb{P})$ is an
$\mathcal{F}$-measurable function.

\paragraph{Meaning of $X_n$ being $\mathcal{F}_n$-measurable.}
The statement means that
\[
\{X_n\in B\}\in\mathcal{F}_n
\qquad
\text{for every }B\in\mathcal{B}(\mathbb{R}).
\]
If $\mathcal{F}_n$ represents the information available at
time $n$, then the value of $X_n$ is determined by that
information.
In particular, every question of the form
``Is $X_n$ in the Borel set $B$?'' can be answered using
the information in $\mathcal{F}_n$.

\paragraph{Example: a sigma-algebra containing partial information.}
Let
\[
\Omega=\{1,2,3,4\},
\qquad
\mathcal{F}=\mathcal{P}(\Omega),
\]
and consider
\[
\mathcal{G}
=
\{\varnothing,\Omega,\{1,2\},\{3,4\}\}.
\]
This information distinguishes the two groups
$\{1,2\}$ and $\{3,4\}$, but does not distinguish outcomes
within either group.

The random variable
\[
U(\omega)=
\begin{cases}
0, & \omega\in\{1,2\},\\
1, & \omega\in\{3,4\}
\end{cases}
\]
is $\mathcal{G}$-measurable.

However,
\[
V(\omega)=
\begin{cases}
0, & \omega\in\{1,3\},\\
1, & \omega\in\{2,4\}
\end{cases}
\]
is not $\mathcal{G}$-measurable, because
\[
\{V=1\}=\{2,4\}\notin\mathcal{G}.
\]
In this finite example, a function is
$\mathcal{G}$-measurable precisely when it is constant
on each of the two groups.

\subsection{Indicators and integrability}

\paragraph{Definition: indicator random variable.}
For an event $A\in\mathcal{F}$, define
\[
\mathbf{1}_A(\omega)=
\begin{cases}
1, & \omega\in A,\\
0, & \omega\notin A.
\end{cases}
\]
Thus,
\[
X\mathbf{1}_A=
\begin{cases}
X, & \text{on }A,\\
0, & \text{on }A^c.
\end{cases}
\]
In particular,
\[
\mathbb{E}[\mathbf{1}_A]=\mathbb{P}(A),
\qquad
\mathbb{E}[X\mathbf{1}_A]
=
\int_A X\,d\mathbb{P}.
\]

\paragraph{Definition: integrability.}
A random variable $X$ is \emph{integrable} if
\[
\mathbb{E}[|X|]<\infty.
\]
If $X$ is integrable, then $X\mathbf{1}_A$ is integrable
for every event $A$, since
\[
|X\mathbf{1}_A|\leq |X|
\quad\Longrightarrow\quad
\mathbb{E}[|X\mathbf{1}_A|]
\leq\mathbb{E}[|X|]<\infty.
\]

\subsection{Filtrations and adapted processes}

\paragraph{Definition: filtration.}
A \emph{filtration} is a sequence
$(\mathcal{F}_n)_{n\geq0}$ of sub-$\sigma$-algebras of
$\mathcal{F}$ such that
\[
\mathcal{F}_n\subseteq\mathcal{F}_{n+1}\subseteq\mathcal{F},
\qquad n\geq0.
\]
The $\sigma$-algebra $\mathcal{F}_n$ represents the information
available at time $n$.
The inclusions express that information accumulates over time.

\paragraph{Definition: stochastic process.}
A discrete-time \emph{stochastic process} is a sequence
\[
X=(X_n)_{n\geq0}
\]
of random variables $X_n:\Omega\to\mathbb{R}$.

\paragraph{Definition: adaptedness.}
The process $X$ is \emph{adapted} to
$(\mathcal{F}_n)_{n\geq0}$ if $X_n$ is
$\mathcal{F}_n$-measurable for every $n\geq0$.

The process is \emph{integrable} if
\[
\mathbb{E}[|X_n|]<\infty
\qquad\text{for every }n\geq0.
\]
This is a condition at each fixed time; it does not require
$\sup_n\mathbb{E}[|X_n|]<\infty$.

\subsection{Natural filtrations}

\paragraph{Definition: natural filtration.}
For a stochastic process $X=(X_n)_{n\geq0}$, its
\emph{natural filtration} is
\[
\mathcal{F}_n^X=\sigma(X_0,X_1,\ldots,X_n),
\qquad n\geq0.
\]
Here $\sigma(X_0,\ldots,X_n)$ denotes the smallest
$\sigma$-algebra making all of $X_0,\ldots,X_n$ measurable.

More explicitly, define the random vector
\[
W_n:\Omega\to\mathbb{R}^{n+1},
\qquad
W_n(\omega)=
\bigl(X_0(\omega),\ldots,X_n(\omega)\bigr).
\]
Then
\[
\mathcal{F}_n^X
=
\left\{
W_n^{-1}(B):
B\in\mathcal{B}(\mathbb{R}^{n+1})
\right\},
\]
or, equivalently,
\[
\mathcal{F}_n^X
=
\left\{
\{\omega\in\Omega:
(X_0(\omega),\ldots,X_n(\omega))\in B\}:
B\in\mathcal{B}(\mathbb{R}^{n+1})
\right\}.
\]

\paragraph{Interpreting the notation.}
For each Borel set $B\subseteq\mathbb{R}^{n+1}$,
\[
\{(X_0,\ldots,X_n)\in B\}
\]
is an event in $\Omega$.
It consists of the outcomes for which the observed vector
belongs to $B$.

For example, when $n=1$, take
\[
B=(-\infty,0]\times(1,\infty).
\]
Then
\[
\{(X_0,X_1)\in B\}
=
\{X_0\leq0\}\cap\{X_1>1\}.
\]
The set $B$ describes a condition on possible observed values;
its preimage is the corresponding event in the sample space.
Allowing $B$ to range over all Borel sets gives the entire
$\sigma$-algebra $\sigma(X_0,X_1)$.

\paragraph{Smallest-filtration property.}
The natural filtration is the smallest filtration to which
$X$ is adapted.

Indeed, suppose $X$ is adapted to another filtration
$(\mathcal{F}_n)_{n\geq0}$.
For $k\leq n$, the random variable $X_k$ is
$\mathcal{F}_k$-measurable and hence
$\mathcal{F}_n$-measurable, since
$\mathcal{F}_k\subseteq\mathcal{F}_n$.
Therefore,
\[
\mathcal{F}_n^X
=
\sigma(X_0,\ldots,X_n)
\subseteq\mathcal{F}_n.
\]

\subsection{Conditional probabilities and expectations}

\subsubsection{Conditioning on an event}

Let $A,B\in\mathcal{F}$ with $\mathbb{P}(B)>0$.
The conditional probability of $A$ given $B$ is
\[
\mathbb{P}(A\mid B)
=
\frac{\mathbb{P}(A\cap B)}{\mathbb{P}(B)}.
\]
For fixed $B$, the map
\[
A\longmapsto\mathbb{P}(A\mid B)
\]
is a probability measure on $(\Omega,\mathcal{F})$.

For an integrable random variable $X$, define
\[
\mathbb{E}[X\mid B]
=
\frac{\mathbb{E}[X\mathbf{1}_B]}{\mathbb{P}(B)}.
\]
This is a number: the average of $X$ conditional on the
event $B$ occurring.

\subsubsection{Conditioning on a sigma-algebra}

\paragraph{Theorem: existence and almost sure uniqueness.}
Let $\mathcal{G}\subseteq\mathcal{F}$ be a $\sigma$-algebra,
and let $X$ be integrable.
There exists an integrable random variable $Y$ satisfying:
\begin{enumerate}
    \item $Y$ is $\mathcal{G}$-measurable;
    \item for every $A\in\mathcal{G}$,
    \[
    \mathbb{E}[Y\mathbf{1}_A]
    =
    \mathbb{E}[X\mathbf{1}_A].
    \]
\end{enumerate}
Moreover, if $Y'$ is another integrable random variable
satisfying these conditions for the same $X$, then
\[
Y=Y'\qquad\text{a.s.}
\]

\paragraph{Definition: conditional expectation.}
Any such $Y$ is called a \emph{version of the conditional
expectation of $X$ given $\mathcal{G}$}.
We write
\[
Y=\mathbb{E}[X\mid\mathcal{G}]
\qquad\text{a.s.}
\]

The two defining conditions have distinct roles:
\begin{itemize}
    \item Measurability requires $Y$ to use only the information
    contained in $\mathcal{G}$.
    \item The integral identities require $Y$ to preserve
    the average of $X$ over every event observable using
    that information:
    \[
    \int_A Y\,d\mathbb{P}
    =
    \int_A X\,d\mathbb{P},
    \qquad A\in\mathcal{G}.
    \]
\end{itemize}

Taking $A=\Omega$ gives
\[
\mathbb{E}[Y]=\mathbb{E}[X].
\]
However, this single equality is insufficient to characterize
conditional expectation: the identity must hold for every
$A\in\mathcal{G}$.

\paragraph{Why the notation includes ``almost surely''.}
The theorem identifies a unique random variable up to
almost sure equality.
Two admissible versions may differ on a set of probability
zero while satisfying exactly the same integral identities.

Thus, writing
\[
Y=\mathbb{E}[X\mid\mathcal{G}]
\qquad\text{a.s.}
\]
names the random variable characterized by the theorem.
It is not an additional assumption on $Y$.

\subsubsection{Monotonicity and the proof of uniqueness}

Suppose $X$ and $X'$ are integrable and
\[
X\leq X'\qquad\text{a.s.}
\]
Let $Y$ and $Y'$ satisfy the defining conditions of
conditional expectation for $X$ and $X'$, respectively,
with respect to the same $\sigma$-algebra $\mathcal{G}$.

Set
\[
A=\{Y>Y'\},
\qquad
Z=(Y-Y')\mathbf{1}_A.
\]
Since $Y$ and $Y'$ are $\mathcal{G}$-measurable,
\[
A\in\mathcal{G}.
\]
Also,
\[
0\leq Z\leq |Y|+|Y'|,
\]
so $Z$ is nonnegative and integrable.

Using the defining integral identities,
\begin{align*}
\mathbb{E}[Z]
&=
\mathbb{E}[(Y-Y')\mathbf{1}_A]\\
&=
\mathbb{E}[Y\mathbf{1}_A]
-
\mathbb{E}[Y'\mathbf{1}_A]\\
&=
\mathbb{E}[X\mathbf{1}_A]
-
\mathbb{E}[X'\mathbf{1}_A]\\
&=
\mathbb{E}[(X-X')\mathbf{1}_A]
\leq0.
\end{align*}
Since $Z\geq0$, we also have $\mathbb{E}[Z]\geq0$.
Hence
\[
\mathbb{E}[Z]=0.
\]
A nonnegative random variable with expectation zero is
zero almost surely, so $Z=0$ a.s.
Consequently,
\[
Y\leq Y'\qquad\text{a.s.}
\]

This proves \emph{monotonicity}:
\[
X\leq X'\quad\text{a.s.}
\quad\Longrightarrow\quad
\mathbb{E}[X\mid\mathcal{G}]
\leq
\mathbb{E}[X'\mid\mathcal{G}]
\quad\text{a.s.}
\]

For uniqueness, take $X'=X$.
The argument gives $Y\leq Y'$ a.s.
Interchanging $Y$ and $Y'$ gives $Y'\leq Y$ a.s.
Therefore,
\[
Y=Y'\qquad\text{a.s.}
\]

\subsubsection{A finite-partition example}

Suppose $A_1,\ldots,A_m$ form a measurable partition
of $\Omega$, with $\mathbb{P}(A_i)>0$, and let
\[
\mathcal{G}=\sigma(A_1,\ldots,A_m).
\]
Then
\[
\mathbb{E}[X\mid\mathcal{G}]
=
\sum_{i=1}^{m}
\frac{\mathbb{E}[X\mathbf{1}_{A_i}]}{\mathbb{P}(A_i)}
\mathbf{1}_{A_i}.
\]
Thus, on each part of the partition, conditional expectation
replaces $X$ by its average over that part.

For the earlier example
\[
\Omega=\{1,2,3,4\},
\qquad
\mathcal{G}
=
\{\varnothing,\Omega,\{1,2\},\{3,4\}\},
\]
suppose all four outcomes are equally likely and $X(\omega)=\omega$.
Then
\[
\mathbb{E}[X\mid\mathcal{G}](\omega)
=
\begin{cases}
\frac{1+2}{2}=\frac32, & \omega\in\{1,2\},\\[4pt]
\frac{3+4}{2}=\frac72, & \omega\in\{3,4\}.
\end{cases}
\]
This is a random variable whose value depends on which
group is observed.

\subsubsection{Conditioning on a random variable}

For a real-valued random variable $W$, define
\[
\sigma(W)
=
\bigl\{
W^{-1}(B):B\in\mathcal{B}(\mathbb{R})
\bigr\}.
\]
Conditioning on $W$ means conditioning on the information
generated by $W$:
\[
\mathbb{E}[X\mid W]
:=
\mathbb{E}[X\mid\sigma(W)].
\]

\subsubsection{Conditional probability given a sigma-algebra}

For an event $A\in\mathcal{F}$, define
\[
\mathbb{P}(A\mid\mathcal{G})
:=
\mathbb{E}[\mathbf{1}_A\mid\mathcal{G}].
\]
This is a $\mathcal{G}$-measurable random variable, with
\[
0\leq\mathbb{P}(A\mid\mathcal{G})\leq1
\qquad\text{a.s.}
\]
Its defining integral identity is
\[
\mathbb{E}\!\left[
\mathbb{P}(A\mid\mathcal{G})\mathbf{1}_B
\right]
=
\mathbb{P}(A\cap B),
\qquad B\in\mathcal{G}.
\]

\subsection{Martingales, submartingales, and supermartingales}

Fix a filtration $(\mathcal{F}_n)_{n\geq0}$.

\paragraph{Definition: martingale.}
A process $X=(X_n)_{n\geq0}$ is a \emph{martingale}
with respect to $(\mathcal{F}_n)_{n\geq0}$ if:
\begin{enumerate}
    \item $X$ is adapted:
    $X_n$ is $\mathcal{F}_n$-measurable for every $n$;
    \item $X$ is integrable:
    $\mathbb{E}[|X_n|]<\infty$ for every $n$;
    \item the \emph{martingale property} holds:
    \[
    \mathbb{E}[X_{n+1}\mid\mathcal{F}_n]
    =
    X_n
    \qquad\text{a.s., for every }n\geq0.
    \]
\end{enumerate}

The martingale property says that, given the information
available at time $n$, the conditional expected value at
the next time is the current value.

\paragraph{Definition: submartingale.}
An adapted, integrable process $X$ is a
\emph{submartingale} if
\[
\mathbb{E}[X_{n+1}\mid\mathcal{F}_n]
\geq X_n
\qquad\text{a.s., for every }n\geq0.
\]

\paragraph{Definition: supermartingale.}
An adapted, integrable process $X$ is a
\emph{supermartingale} if
\[
\mathbb{E}[X_{n+1}\mid\mathcal{F}_n]
\leq X_n
\qquad\text{a.s., for every }n\geq0.
\]

Thus, the conditional expected increment is respectively
zero, nonnegative, or nonpositive:
\[
\mathbb{E}[X_{n+1}-X_n\mid\mathcal{F}_n]
\begin{cases}
=0, & \text{for a martingale},\\
\geq0, & \text{for a submartingale},\\
\leq0, & \text{for a supermartingale}.
\end{cases}
\]

\subsection{Equivalent integral characterizations}

Suppose $X$ is adapted and integrable.
Then:
\begin{align*}
X\text{ is a martingale}
\quad&\Longleftrightarrow\quad
\mathbb{E}[X_{n+1}\mathbf{1}_A]
=
\mathbb{E}[X_n\mathbf{1}_A],\\
X\text{ is a submartingale}
\quad&\Longleftrightarrow\quad
\mathbb{E}[X_{n+1}\mathbf{1}_A]
\geq
\mathbb{E}[X_n\mathbf{1}_A],\\
X\text{ is a supermartingale}
\quad&\Longleftrightarrow\quad
\mathbb{E}[X_{n+1}\mathbf{1}_A]
\leq
\mathbb{E}[X_n\mathbf{1}_A],
\end{align*}
where each condition on the right must hold for
\emph{every} $n\geq0$ and \emph{every} $A\in\mathcal{F}_n$.

\paragraph{Proof of the martingale characterization.}
Fix $n$ and let
\[
Y=\mathbb{E}[X_{n+1}\mid\mathcal{F}_n].
\]
By definition,
\[
\mathbb{E}[Y\mathbf{1}_A]
=
\mathbb{E}[X_{n+1}\mathbf{1}_A],
\qquad A\in\mathcal{F}_n.
\]
If $Y=X_n$ a.s., the required integral equality follows.

Conversely, suppose
\[
\mathbb{E}[X_{n+1}\mathbf{1}_A]
=
\mathbb{E}[X_n\mathbf{1}_A],
\qquad A\in\mathcal{F}_n.
\]
Since $X_n$ is integrable and $\mathcal{F}_n$-measurable,
it satisfies the defining conditions for the conditional
expectation of $X_{n+1}$ given $\mathcal{F}_n$.
Almost sure uniqueness therefore gives
\[
X_n=\mathbb{E}[X_{n+1}\mid\mathcal{F}_n]
\qquad\text{a.s.}
\]

\paragraph{Why the inequality characterizations also hold.}
Again let $Y=\mathbb{E}[X_{n+1}\mid\mathcal{F}_n]$.
If $Y\geq X_n$ a.s., then
\[
\mathbb{E}[(Y-X_n)\mathbf{1}_A]\geq0
\qquad\text{for every }A\in\mathcal{F}_n.
\]

Conversely, suppose these integral inequalities hold.
The event
\[
A=\{Y<X_n\}
\]
belongs to $\mathcal{F}_n$.
On this event, $(Y-X_n)\mathbf{1}_A\leq0$, while the
assumed inequality gives
\[
\mathbb{E}[(Y-X_n)\mathbf{1}_A]\geq0.
\]
Hence this expectation is zero, so
\[
(X_n-Y)\mathbf{1}_{\{Y<X_n\}}=0
\qquad\text{a.s.}
\]
Therefore $Y\geq X_n$ a.s.
The supermartingale case follows by reversing the inequalities.

\paragraph{Consequence for ordinary expectations.}
Taking $A=\Omega$ shows that
\[
\mathbb{E}[X_{n+1}]
\begin{cases}
=\mathbb{E}[X_n], & \text{for a martingale},\\
\geq\mathbb{E}[X_n], & \text{for a submartingale},\\
\leq\mathbb{E}[X_n], & \text{for a supermartingale}.
\end{cases}
\]
These conditions on ordinary expectations alone are not
sufficient: the integral conditions must hold on every
event in $\mathcal{F}_n$.

\subsection{A martingale under its natural filtration}

\paragraph{Proposition.}
If $X$ is a martingale with respect to
$(\mathcal{F}_n)_{n\geq0}$, then it is also a martingale
with respect to its natural filtration
$(\mathcal{F}_n^X)_{n\geq0}$.

\paragraph{Proof.}
The process $X$ is adapted to its natural filtration,
and each $X_n$ remains integrable.
Moreover,
\[
\mathcal{F}_n^X\subseteq\mathcal{F}_n.
\]
Thus, for every $A\in\mathcal{F}_n^X$, the original
martingale property gives
\[
\mathbb{E}[X_{n+1}\mathbf{1}_A]
=
\mathbb{E}[X_n\mathbf{1}_A].
\]
The integral characterization proves that $X$ is a
martingale with respect to $(\mathcal{F}_n^X)_{n\geq0}$.

\subsection{Examples from independent random variables}

Let $X_1,X_2,\ldots$ be independent integrable random variables.
Set
\[
\mathcal{F}_0=\{\varnothing,\Omega\},
\qquad
\mathcal{F}_n=\sigma(X_1,\ldots,X_n),
\qquad n\geq1.
\]
Then $X_{n+1}$ is independent of $\mathcal{F}_n$.

Define
\[
S_0=0,
\qquad
S_n=\sum_{i=1}^{n}X_i,
\]
and
\[
Z_0=1,
\qquad
Z_n=\prod_{i=1}^{n}X_i.
\]

\subsubsection{Sums of independent mean-zero random variables}

Suppose
\[
\mathbb{E}[X_n]=0
\qquad\text{for every }n\geq1.
\]
Then $(S_n)_{n\geq0}$ is a martingale with respect to
$(\mathcal{F}_n)_{n\geq0}$.

\paragraph{Proof.}
Each $S_n$ is a measurable function of $X_1,\ldots,X_n$,
so $S_n$ is $\mathcal{F}_n$-measurable.
Also,
\[
\mathbb{E}[|S_n|]
\leq
\sum_{i=1}^{n}\mathbb{E}[|X_i|]
<\infty.
\]

For $A\in\mathcal{F}_n$, independence gives
\[
\mathbb{E}[X_{n+1}\mathbf{1}_A]
=
\mathbb{E}[X_{n+1}]\,\mathbb{P}(A)
=0.
\]
Therefore,
\begin{align*}
\mathbb{E}[S_{n+1}\mathbf{1}_A]
&=
\mathbb{E}[(S_n+X_{n+1})\mathbf{1}_A]\\
&=
\mathbb{E}[S_n\mathbf{1}_A]
+
\mathbb{E}[X_{n+1}\mathbf{1}_A]\\
&=
\mathbb{E}[S_n\mathbf{1}_A].
\end{align*}
The integral characterization shows that $S$ is a martingale.

\subsubsection{Products of independent nonnegative mean-one variables}

Suppose
\[
X_n\geq0\quad\text{a.s.},
\qquad
\mathbb{E}[X_n]=1
\qquad\text{for every }n\geq1.
\]
Then $(Z_n)_{n\geq0}$ is a martingale with respect to
$(\mathcal{F}_n)_{n\geq0}$.

\paragraph{Proof.}
Each $Z_n$ is a measurable function of $X_1,\ldots,X_n$,
so $Z_n$ is $\mathcal{F}_n$-measurable.

Nonnegativity and independence give
\[
\mathbb{E}[|Z_n|]
=
\mathbb{E}[Z_n]
=
\prod_{i=1}^{n}\mathbb{E}[X_i]
=1.
\]
Thus, $Z_n$ is integrable.

For $A\in\mathcal{F}_n$, the random variable
$Z_n\mathbf{1}_A$ is $\mathcal{F}_n$-measurable and
integrable.
It is therefore independent of $X_{n+1}$.
Consequently,
\begin{align*}
\mathbb{E}[Z_{n+1}\mathbf{1}_A]
&=
\mathbb{E}[X_{n+1}Z_n\mathbf{1}_A]\\
&=
\mathbb{E}[X_{n+1}]
\,\mathbb{E}[Z_n\mathbf{1}_A]\\
&=
\mathbb{E}[Z_n\mathbf{1}_A].
\end{align*}
The integral characterization shows that $Z$ is a martingale.
\end{document}
